% Generated draft: requires mathematical and source review.
\documentclass[11pt]{article}
\usepackage[margin=25mm]{geometry}
\usepackage{fontspec}
\setmainfont{Latin Modern Roman}
\usepackage{amsmath,amssymb,mathtools,unicode-math}
\setmathfont{Latin Modern Math}
\usepackage{enumitem,needspace,xurl,hyperref}
\setlength{\parindent}{0pt}
\setlength{\parskip}{5pt}
\setlength{\emergencystretch}{3em}
\newcommand{\N}{\mathbb{N}}
\newcommand{\Z}{\mathbb{Z}}
\newcommand{\Q}{\mathbb{Q}}
\newcommand{\R}{\mathbb{R}}
\newcommand{\C}{\mathbb{C}}
\newcommand{\F}{\mathbb{F}}
\newcommand{\A}{\mathbb{A}}
\newcommand{\PP}{\mathbb{P}}
\newcommand{\E}{\mathbb{E}}
\newcommand{\eps}{\varepsilon}
\newcommand{\dd}{\,\mathrm{d}}
\begin{document}
\section*{UnsolvedMath ID: 30004305}
\textbf{Source problem code:} catalogue entry
\subsection{Definitions and mathematical statement}
For notation, let \(X\) be a set and let \(r\colon X\times X\to X\times X\) satisfy the braid form of the set-theoretic Yang--Baxter equation
\[
(r\times\operatorname{id}_X)(\operatorname{id}_X\times r)(r\times\operatorname{id}_X)
=
(\operatorname{id}_X\times r)(r\times\operatorname{id}_X)(\operatorname{id}_X\times r)
\quad\text{on }X^3.
\]
Write \(r(x,y)=(\lambda_x(y),\rho_y(x))\). The associated structure monoid is
\[
M(X,r)=\left\langle X\ \middle|\ xy=\lambda_x(y)\rho_y(x)\text{ for all }x,y\in X\right\rangle.
\]
The notation \(kM(X,r)\) denotes its monoid algebra: finite formal \(k\)-linear combinations of monoid elements, with multiplication extending the monoid multiplication. The supplied record does not specify the coefficient assumptions on \(k\), including whether it must be a field. It also does not specify whether \(X\) is finite or whether \(r\) must be bijective, involutive, or nondegenerate. These conditions are therefore not imposed here.

Put \(A=kM(X,r)\). In the usual ring-theoretic sense, a prime ideal \(P\) of \(A\) is a proper two-sided ideal such that, for all two-sided ideals \(I,J\) of \(A\),
\[
IJ\subseteq P\quad\Longrightarrow\quad I\subseteq P\ \text{or}\ J\subseteq P.
\]
Its prime quotient is \(A/P\).

\textbf{Source question.} For those prime ideals \(P\) for which \(kM(X,r)/P\) arises from a Yang--Baxter solution related to \((X,r)\), determine how the properties of those related solutions, as \(P\) varies, reflect properties of \((X,r)\).

\textbf{Unresolved formulation.} The record supplies neither a definition of a related solution nor a construction associating such a solution to a prime quotient. In particular, it does not state whether the quotient must itself be isomorphic to a structure algebra, whether an associated solution is unique, or which properties and implications are sought. Thus the question above preserves the supplied research target but cannot be made into a uniquely specified quantified conjecture from this record alone.
\subsection{Short English statement}
Starting with the structure algebra of a Yang--Baxter solution \((X,r)\), consider its prime quotients that arise from related solutions. The question asks what the properties of these solutions, taken over the relevant prime ideals, reveal about \((X,r)\). The meaning of the association and the properties to be compared remain unspecified.
\subsection{Sources}
\begin{itemize}
\item Local registry: \texttt{data/unsolvedmath/problems.json}, with supplementary records in \texttt{catalog/problems.json}; supplied complete source record for permanent UnsolvedMath ID \(30004305\), source problem code \texttt{catalogue entry}, entitled \emph{Related Solutions from Prime Structure-Algebra Quotients}. The supplied set is \emph{Oberwolfach Reports Open Problems}. No individual proposer is identified.
\item Supplied report attribution: \emph{Mini-Workshop: Algebraic Tools for Solving the Yang--Baxter Equation (2020)}, Oberwolfach Reports, pp.\ 3207--3242, DOI \texttt{10.4171/owr/2019/51}. \url{https://doi.org/10.4171/owr/2019/51}. The registry separately gives proposed year 2019. The background describes Oberwolfach Reports as documenting mathematical workshops and their open-problem discussions.
\item Supplied literature reference: Colazzo--Jespers--Van Antwerpen--Verwimp, ``structure monoid and algebra.'' \url{https://arxiv.org/abs/1812.02026}.
\item Supplied literature reference: ``2026 structure-algebra results.'' \url{https://doi.org/10.1016/j.jalgebra.2026.04.044}.
\item Supplied literature reference: ``original OWR report.'' \url{https://doi.org/10.4171/owr/2019/51}.
\item Literature assessment checked at \texttt{2026-08-22}: this is supplied dated curation, not independently verified here. The supplied status is \emph{partially solved}. The assessment states that semiprime structure algebras admit strong embedding and representation results in nondegenerate settings, that recent work treats further finite degenerate cases, and that the requested matrix representations and their effect on solutions are not general theorems. The background identifies this assessment as dated literature triage derived from the supplied OWR curation and requests independent verification of time-sensitive status and cited claims.
\end{itemize}
\end{document}
